\documentclass[12pt,a4paper]{article}

\input{/home/thales/documents/profia/meus_materiais/configuracao_latek/config_profia_1.tex}

\begin{document}

\cabecalhoRoteiro
    {/home/thales/documents/profia/imagens/Logo Profia.png} % Logo 1
    {/home/thales/documents/profia/imagens/Logo Horizontal Sem Fundo.png} % Logo 2
    {Profia - Escola de Programação e Tecnologia}
    {Orientação a Objetos}
    {Introdução ao desenvolvimento Web}
    {Miniprojeto - Blog de Receitas}

\noindent\textbf{Objetivo:} Construir um site de receitas completo integrando HTML Semântico, Box Model e Flexbox passo a passo.

\vspace{0.5cm}

\section*{Parte 1: Estruturando com o HTML Semântico}
Nesta primeira etapa, vamos criar o esqueleto da nossa página bloco a bloco. Focaremos apenas no significado do conteúdo, a beleza vem depois!

% Passo 1
\vspace{0.25cm}
\noindent\textbf{Passo 1: Estrutura Básica}\\
Crie o seu arquivo HTML e defina a base da página. Lembre-se de configurar o idioma para Português, definir o título da aba e fazer a ligação com o CSS.

\begin{codehtml}
<!DOCTYPE html>
<html lang="pt-br">
    <head>
        <meta charset="UTF-8">
        <meta name="viewport" content="width=device-width, initial-scale=1.0">
        <title>Site de Receitas</title>
        <link rel="stylesheet" href="style.css">
    </head>
    <body>
        <!-- Todo o nosso site ficará aqui dentro -->
    </body>
</html>
\end{codehtml}

% Passo 2
\noindent\textbf{Passo 2: Cabeçalho da Página (Header)}\\
Dentro do \verb|<body>|, crie o topo do site com o título principal e uma área de navegação contendo os links.

\begin{codehtml}
        <header>
            <h1>Site de Receitas</h1>
            <nav>
                <ul>
                    <li><a href="\">INÍCIO</a></li>
                    <li><a href="\">DOCES</a></li>
                    <li><a href="\">SALGADOS</a></li>
                </ul>
            </nav>
        </header>
\end{codehtml}

% Passo 3
\noindent\textbf{Passo 3: Iniciando a Receita e a Imagem}\\
Abaixo do cabeçalho, abra o contêiner principal (\verb|<main>|) e o artigo da receita. Adicione o título do prato e a imagem em destaque.

\begin{codehtml}
        <main>
            <article>
                <section>
                    <h2>Pudim de leite</h2>
                    <figure>
                        <img src="imagens/pudim.jpg" alt="Imagem de um lindo pudim">
                        <figcaption>A sobremesa mais amada do brasil</figcaption>
                    </figure>

                <!- Restante do código aqui... -->
                </section>
            </article>
        </main>
\end{codehtml}

% Passo 4
\noindent\textbf{Passo 4: Descrição e Ingredientes (Listas com Marcadores)}\\
Ainda dentro da mesma seção, adicione o texto descritivo. Em seguida, use listas não ordenadas (\verb|<ul>|) para os ingredientes, pois a ordem em que os separamos não altera a receita.

\begin{codehtml}
                    <p>Esta é a melhor e mais fácil receita que você vai encontrar. O segredo é misturar muito bem mas com calma (para não deixar muito ar entrar na receita).</p>

                    <h3>Ingredientes da Massa:</h3>
                    <ul>
                        <li>1 lata de leite condensado</li>
                        <li>1 lata de leite (mesma medida)</li>
                        <li>3 ovos inteiros</li>
                    </ul>

                    <h3>Ingredientes da Calda:</h3>
                    <ul>
                        <li>1 xícara (chá) de açúcar</li>
                        <li>1/2 xícara de água</li>
                    </ul>
\end{codehtml}

% Passo 5
\noindent\textbf{Passo 5: Modo de Preparo (Listas Numeradas)}\\
Para o preparo, a ordem dos passos é fundamental! Por isso, usaremos listas ordenadas (\verb|<ol>|). Feche a \verb|<section>| ao final.

\begin{codehtml}
                    <h3>Modo de Preparo:</h3>

                    <h4>Calda:</h4>
                    <ol>
                        <li>Derreta o açúcar na panela até ficar moreno, acrescente a água e deixe engrossar.</li>
                        <li>Coloque em uma forma redonda e reserve.</li>
                    </ol>

                    <h4>Massa:</h4>
                    <ol>
                        <li>Primeiro, bata bem os ovos no liquidificador.</li>
                        <li>Acrescente o leite condensado e o leite, e bata novamente.</li>
                        <li>Despeje a massa do pudim na forma, por cima da calda.</li>
                        <li>Asse em forno médio por 45 minutos (em banho-maria).</li>
                        <li>Espete um garfo para ver se está bem assado.</li>
                        <li>Deixe esfriar e desenforme.</li>
                    </ol>
                </section>
\end{codehtml}

% Passo 6
\noindent\textbf{Passo 6: O Formulário de Comentários}\\
Inicie uma nova seção para os leitores interagirem. Crie um formulário com os campos de Nome, E-mail, uma caixa de texto grande (\verb|<textarea>|) e o botão de enviar.

\begin{codehtml}
                <section class="comentarios">
                    <form>
                        <h4>Deixe seu comentário: </h4>
                        <label for="nome-formulario">Nome:</label>
                        <input type="text" id="nome-formulario">

                        <label for="email-formulario">E-mail:</label>
                        <input type="email" id="email-formulario">

                        <label for="comentario-formulario">Comentário:</label>
                        <textarea
                            id="comentario-formulario"
                            placeholder="O que você achou da receita?"
                            rows="4"
                            cols="40"
                            minlength="1"
                            wrap="soft"
                            required>
                            Digite aqui seu comentário!
                        </textarea>

                        <button>Enviar Comentário</button>
                    </form>
\end{codehtml}

% Passo 7
\noindent\textbf{Passo 7: Exibindo um Comentário}\\
Logo abaixo do formulário, crie um bloco (\verb|<div>|) simulando um comentário que já foi publicado no site. Feche a seção e o artigo.

\begin{codehtml}
                    <h3>Comentários dos leitores:</h3>
                    <div class="comentario-leitor">
                        <p><strong>Maria:</strong> Fiz ontem e ficou maravilhoso!</p>
                    </div>
                </section>
            </article>
\end{codehtml}

% Passo 8
\noindent\textbf{Passo 8: Barra Lateral (Aside)}\\
Fora do artigo (mas dentro do \verb|<main>|), adicione uma barra lateral recomendando outros pratos aos visitantes.

\begin{codehtml}
            <aside>
                <h2>Veja Também</h2>
                <ul>
                    <li><a href="">Bolo de cenoura com chocolate</a></li>
                    <li><a href="">Torta de frango cremosa</a></li>
                    <li><a href="">Pão de queijo de frigideira</a></li>
                </ul>
            </aside>
        </main>
\end{codehtml}

% Passo 9
\noindent\textbf{Passo 9: Rodapé (Footer)}\\
Para finalizar o HTML, insira o rodapé da página com os direitos autorais.

\begin{codehtml}
        <footer>
            <p>&copy; Blog Site de Receitas. Todos os direito reservados - 2026</p>
        </footer>
\end{codehtml}

\newpage

\section*{Parte 2: A Mágica do CSS e Flexbox}
Agora que a estrutura está sólida, vamos usar o arquivo \verb|style.css| para organizar o layout e pintar o nosso site.

% CSS Passo 1
\vspace{0.25cm}
\noindent\textbf{Passo 1: Reset e Configuração Global}\\
Zere as margens universais e defina a fonte e a cor de fundo padrão para todo o site.

\begin{codecss}
* {
    margin: 0;
    padding: 0;
    box-sizing: border-box;
}

body{
    font-family: Arial, Helvetica, sans-serif;
    background-color: #f4f4f9;
    color: #333333;
}
\end{codecss}

% CSS Passo 2
\noindent\textbf{Passo 2: O Cabeçalho (Header)}\\
Pinte o fundo do cabeçalho de uma cor vibrante, centralize o título e dê um espaçamento interno (\verb|padding|).

\begin{codecss}
header {
    background-color: #ff6b6b;
    color: white;
    text-align: center;
    padding: 10px 15px;
}

header h1 {
    color:#333
}
\end{codecss}

% CSS Passo 3
\noindent\textbf{Passo 3: O Menu (Flexbox)}\\
Use o Flexbox na lista (\verb|<ul>|) do menu para colocar os itens um ao lado do outro, centralizados. Remova as bolinhas da lista e os sublinhados dos links.

\begin{codecss}
nav ul {
    display: flex;
    justify-content: center;
    list-style: none;
    margin-top: 15px;
}

nav li {
    margin: 0 15px;
}

nav a {
    color: white;
    text-decoration: none;
    font-weight: bold;
}
\end{codecss}

% CSS Passo 4
\noindent\textbf{Passo 4: O Layout Principal (Duas Colunas)}\\
Vamos criar as colunas do site! Transforme o \verb|<main>| em um contêiner Flex, fazendo o \verb|<article>| ocupar mais espaço (\verb|flex: 3|) e a barra lateral \verb|<aside>| ser mais estreita (\verb|flex: 1|).

\begin{codecss}
main {
    display: flex;
    max-width: 1000px;
    margin: 30px auto; /* Centraliza o site na tela */
    gap: 20px;         /* Espaço entre as colunas */
    padding: 0 20px;
}

article {
    flex: 3;
    background-color: white;
    padding: 20px;
    border-radius: 8px
}

aside {
    flex: 1;
    background-color: #fff0f0;
    padding: 20px;
    border-radius: 8px;
}
\end{codecss}

% CSS Passo 5
\noindent\textbf{Passo 5: A Imagem da Receita}\\
Garanta que a imagem nunca ultrapasse o tamanho da tela, forçando-a a ter 100\% da largura disponível na coluna. Arredonde as bordas.

\begin{codecss}
figure img {
    width: 100%;
    border-radius: 8px;
    margin-top: 15px;
}

figcaption {
    text-align: center;
    font-style: italic;
    color: #777;
    margin-bottom: 20px;
}
\end{codecss}

% CSS Passo 6
\noindent\textbf{Passo 6: Tipografia (Títulos, Textos e Listas)}\\
Adicione respiros (margens) nos subtítulos para que não fiquem colados no texto. Aumente o espaço entre as linhas (\verb|line-height|) para melhorar a leitura.

\begin{codecss}
article h3 {
    margin-top: 25px;
    margin-bottom: 15px;
}

article h4 {
    margin-top: 15px;
    margin-bottom: 10px;
}

p, ul, ol {
    margin-bottom: 15px;
    line-height: 1.6;
}

article li {
    margin-left: 30px;
}

aside ul {
    list-style-type: square;
    margin-left: 20px;
    margin-top: 15px;
}

aside a {
    color: #ff6b6b;
    text-decoration: none;
}
\end{codecss}

% CSS Passo 7
\noindent\textbf{Passo 7: O Formulário}\\
Use Flexbox novamente no \verb|<form>|, mas agora com a direção de \textbf{coluna}, para que os campos caiam um abaixo do outro automaticamente.

\begin{codecss}
form {
    display: flex;
    flex-direction: column;
    gap: 10px;
    margin-bottom: 20px;
}

input, textarea {
    padding: 10px;
    border: 1px solid #ccc;
    border-radius: 4px;
    font-family: inherit;
    resize: none; /* Impede o usuário de redimensionar a caixa */
}
\end{codecss}

% CSS Passo 8
\noindent\textbf{Passo 8: O Botão e a Área de Comentários}\\
Destaque o botão com uma cor forte e mude o cursor quando o mouse passar por cima (\verb|hover|). Dê uma cor de fundo sutil para a caixa de comentários postados.

\begin{codecss}
button {
    background-color: #ff6b6b;
    color: white;
    border: none;
    padding: 12px;
    font-weight: bold;
    cursor: pointer;
}

button:hover{
    background-color: #e55a5a;
}

.comentarios{
    margin-top: 40px;
    border-top: 2px solid #eee;
    padding-top: 20px;
}

.comentario-leitor {
    background-color: #fae7e7;
    padding: 15px;
    border-radius: 4px;
    margin-bottom: 20px;
    margin-top: 10px;
}
\end{codecss}

% CSS Passo 9
\noindent\textbf{Passo 9: O Rodapé}\\
Para fechar com estilo, dê um fundo escuro ao rodapé e centralize o texto.

\begin{codecss}
footer {
    text-align: center;
    background-color: #333;
    color: white;
    padding: 20px;
    margin-top: 40px;
}
\end{codecss}

\end{document}
